% ================================================================
%  Lectura del pipeline con patrones EIP
% ================================================================
\section{Lectura del pipeline con patrones EIP}

La Tabla~\ref{tab:eip} nombra cada etapa implementada con su patrón de Hohpe \& Woolf.

\begin{longtable}{L{3.2cm} L{3.8cm} L{3.8cm} L{4.2cm}}
  \caption{Etapas del pipeline y patrones EIP}\label{tab:eip}\\
  \tblhead Etapa & \color{white}\bfseries Implementación & \color{white}\bfseries Patrón & \color{white}\bfseries Por qué\\
  \endfirsthead
  \multicolumn{4}{l}{\small\textit{Tabla~\ref{tab:eip} -- continuación}}\\[2pt]
  \tblhead Etapa & \color{white}\bfseries Implementación & \color{white}\bfseries Patrón & \color{white}\bfseries Por qué\\
  \endhead\bottomrule\endlastfoot
  \multicolumn{4}{l}{\textbf{Ingesta}}\\
  Lectura de los PDF & \texttt{PdfRepository} & Channel Adapter & Conecta el archivo con el flujo\\
  Estructura común & \texttt{SectionParser}, \texttt{ExtractedDocument} & Normalizer, Canonical Data Model & Lleva distintos tipos de documento a un solo modelo\\
  Verificación & \texttt{ExtractionVerifier} & Message Filter & Deja pasar solo los documentos correctos\\
  Partición & \texttt{Partitioner} & Splitter & Un documento produce varios fragmentos\\
  Vigencia y vectores & Copia del registro, \texttt{Embedding\-Repository}, \texttt{Lexical\-Tokenizer} & Content Enricher & Agrega al fragmento datos de otras fuentes\\
  Marca de inyección & \texttt{InjectionDetector} & Content Enricher & Agrega una marca, sin decidir la entrega\\
  Carga idempotente & \texttt{QdrantRepository}, identificador determinista & Idempotent Receiver & Reingerir no duplica\\
  Cambio de vigencia & \texttt{--sync-validity} & Event-Driven Consumer (diseño); hoy, ejecución manual & En producción reaccionaría a un evento\\
  \multicolumn{4}{l}{\textbf{Consulta}}\\
  Herramientas MCP & \texttt{DocumentTools} & Service Activator, Request-Reply & El protocolo invoca el servicio y espera su respuesta\\
  Documento excluido, rechazo o respuesta & \texttt{ResponseComposer} & Content-Based Router & El contenido define la ruta\\
  Filtro de vigencia & \texttt{\_VALID\_ONLY} & Selective Consumer & Solo se consideran los fragmentos vigentes\\
  Ramas semántica y léxica & Consulta híbrida & Scatter-Gather & Una consulta va a dos ramas, y se reúnen sus resultados\\
  Fusión RRF & Qdrant, \texttt{stable\_order} & Aggregator & Combina los resultados en una sola lista\\
  Retención de contenido & \texttt{\_content} & Content Filter & Quita el texto y conserva la cita\\
  Sección padre & \texttt{\_parent\_context} & Content Enricher & Agrega la sección a la respuesta\\
  Cita & Identificador \texttt{code@version\#section} & Correlation Identifier & Une la respuesta con su fragmento de origen\\
\end{longtable}

\subsection{Dónde la analogía deja de sostenerse}

\begin{itemize}
  \item No hay mensajes. El pipeline es un proceso por lotes y llamadas síncronas dentro de un programa. No hay canales, colas ni consumidores independientes, de modo que no tiene las garantías que esos patrones suponen, como entrega garantizada o desacoplamiento temporal.
  \item La decisión es probabilística. Un Content-Based Router clásico evalúa una condición exacta. Con preguntas en lenguaje natural, la ruta depende de una similitud y un umbral calibrado: el mismo tipo de mensaje puede tomar otra ruta con otras palabras. Por eso aparecen respuestas falsas en temas cercanos.
  \item La transformación no es reversible. Un Message Translator conserva el significado. Un embedding pierde información y no permite reconstruir el texto, ni explicar por qué dos textos se parecen.
  \item La fusión no correlaciona. Un Aggregator reúne partes de un mismo mensaje hasta completarlo. RRF mezcla dos clasificaciones por posición, y no sabe si el resultado está completo o es correcto.
  \item El Filter no juzga significado. El filtro de vigencia es exacto; la pertinencia no lo es. Que un fragmento pase el umbral no significa que responda la pregunta (sección~\ref{sec:citas-erradas}).
  \item No existe un Invalid Message Channel. Un documento rechazado solo queda en el log, sin canal para reprocesarlo.
\end{itemize}

El vocabulario sirve para las etapas de datos e integración, donde la analogía es útil: partición, enriquecimiento, idempotencia. Deja de servir donde interviene la similitud, que es precisamente donde está la mayoría de los fallos medidos.
